\subsection{Evaluación de latencia en sistemas robóticos distribuidos}
\label{sec:evaluacion_de_latencia_en_sistemas_roboticos_distribuidos}

En el modelo \textit{publish--subscribe} los mensajes son transmitidos a través de distintos niveles del \textit{middleware} y de la pila de red. Este proceso introduce diferentes fuentes de latencia que pueden afectar al comportamiento global de la aplicación robótica. Entre estas fuentes se incluyen el tiempo de serialización de los mensajes, la gestión de colas de publicación y suscripción, la planificación de procesos en el sistema operativo y el retardo introducido por la red.

Diversos trabajos han analizado el impacto de estos factores en el rendimiento de aplicaciones robóticas. En particular, Nishimura et al. proponen la herramienta \gls{raplet} para analizar la latencia de comunicación entre nodos en aplicaciones basadas en \gls{ros} \cite{nishimura2021raplet}. Esta herramienta permite descomponer la latencia total en distintos componentes asociados al procesamiento de los nodos, la planificación de hilos en el sistema operativo y la transmisión de datos a través de la pila de red. Los resultados experimentales muestran que la latencia de comunicación puede variar significativamente en función de factores como el tamaño de los mensajes, la carga computacional del sistema o la configuración del \textit{hardware} utilizado. En aplicaciones robóticas complejas, en las que múltiples nodos intercambian grandes volúmenes de información sensorial, estos retardos pueden acumularse y afectar al comportamiento temporal del sistema.

Estudios recientes han analizado el comportamiento de este tipo de arquitecturas en entornos experimentales que combinan robots móviles, redes de comunicaciones avanzadas y plataformas de computación distribuida. Por ejemplo, en plataformas robóticas conectadas a redes móviles de nueva generación, la latencia total del sistema puede estar influenciada por múltiples componentes, incluyendo la transmisión de vídeo, el procesamiento de algoritmos de percepción y el envío de comandos de control al robot \cite{Lozano2025DigitalTwin6G}.

En consecuencia, la evaluación de la latencia en sistemas robóticos distribuidos constituye un elemento clave para el diseño y la validación de arquitecturas robóticas. El análisis de los distintos componentes que contribuyen al retardo total permite identificar cuellos de botella en el sistema y optimizar la distribución de funciones entre los distintos elementos de la arquitectura.